\documentclass[10pt,a4paper]{amsart}
\usepackage[margin=19mm]{geometry}
\usepackage{amsmath,amssymb,mathtools}
\usepackage[scr=rsfs,cal=euler]{mathalfa}
\usepackage{xfrac}
\usepackage[unicode,hidelinks,bookmarks=false]{hyperref}
\newcommand{\ie}{i.e.\ }
\newcommand{\cf}{cf.\ }
\newcommand{\resp}{respectively\ }
\newcommand{\Subsection}{\subsection}
\providecommand{\nobreakdash}{\nobreak\hbox{-}}
\setlength{\emergencystretch}{2em}
\multlinegap0pt
\usepackage{amssymb}
\usepackage{mathtools}
\usepackage{xcolor}
\usepackage{braket}
\newtheorem{proposition}{Proposition}
\title{Exact Solutions for Compactly Supported Parabolic and Landau Quantum Barriers}
\author{Peter Collas, David Klein}
\date{Source: arXiv:2601.06369v3 (CC BY 4.0)}
\begin{document}
\maketitle

\noindent\emph{Source and licence.} Exact Solutions for Compactly Supported Parabolic and Landau Quantum Barriers. Version arXiv:2601.06369v3 (CC BY 4.0), distributed by arXiv under the Creative Commons Attribution 4.0 International licence, http://creativecommons.org/licenses/by/4.0/, as recorded in the arXiv licence metadata for this exact version and shown on the abstract page https://arxiv.org/abs/2601.06369v3. Full licence notice: \texttt{licenses/arxiv-2601.06369-CC-BY-4.0.txt}.\par\smallskip

\noindent\textit{Editorial scope.} Counted unit: the article's proposition 1 (source0.tex lines 629--659). The article's complete original proof of it is delivered with it (source0.tex lines 661--749, one \texttt{proof} environment of 89 lines). No mathematics is written, repaired or re-derived: the statement and the proof are the article's own lines, in the article's own notation, and no retained line is substituted or re-flowed.  No further source-local context is needed: every cross-reference of every retained line resolves inside the two delivered spans.  Nothing was omitted from any retained span, so the package carries no omission record.  No retained line contains a citation, so the wrapper carries no bibliography and no bibliography entry.  No retained line contains a figure environment, an image-inclusion command or drawing code, so no image is delivered and the package declares no asset dependency.  Editorial formatting only, in addition: an AMS \texttt{amsart} preamble that restates the article's own declarations and macros used by the retained lines, namely the environments \texttt{proposition} on separate counters, together with the standard packages those lines need.  The retained environments print in this document as Proposition 1 in order of appearance, whereas the article prints the same items with the numbering of its own full document.  The complete native source of the article as distributed in the arXiv e-print for this exact version is bundled unchanged as \texttt{sources/192406/source0.tex}.
\begin{proposition}
Let $\sigma>0$, assume that $\lambda$ is real and that $V(x)$ is a solution of the differential equation
\begin{equation}
\frac{d^2 V}{dx^2}+ (\sigma x^2 - \lambda)V=0.\label{1}
\end{equation}

Then the function 
\begin{equation}
w(z):=V\left(\frac{z}{\sqrt{2}\,\sigma^{(1/4)}}\right),\label{2}
\end{equation}

is a solution to
\begin{equation}
\frac{d^2 w}{dz^2}+\left(\frac{z^2}{4}-\frac{\lambda}{2\sqrt{\sigma}}\right)\!w(z)=0.\label{3}
\end{equation}

Conversely, if $w(z)$ is a solution to
\begin{equation}
\frac{d^2 w}{dz^2}+\left(\frac{z^2}{4}-a\right)w(z)=0,\label{4}
\end{equation}

where the parameter $a$ is real, then the function
\begin{equation}
V(x)=w(\sqrt{2}\,\sigma^{(1/4)}x),\label{5}
\end{equation}

is a solution to
\begin{equation}
\frac{d^2 V}{dx^2}+ (\sigma x^2 - 2a\sqrt{\sigma})V=0.\label{6}
\end{equation}
\end{proposition}

\begin{proof}
Let $x$ be a function of $z$ given by
\begin{equation}
x(z)= \frac{z}{\sqrt{2}\,\sigma^{(1/4)}}\label{7}
\end{equation}

and define,
\begin{equation}
w(z):=V(x(z))= V\left(\frac{z}{\sqrt{2}\,\sigma^{(1/4)}}\right),\label{8}
\end{equation}

So we may regard $V=V(x(z))$ as a function of $z$ and the same is true of the second derivative of $V$ with respect to $x$. 

then by the chain rule,
\begin{equation}
\frac{dw}{dz}=\frac{dV}{dx}\frac{dx}{dz}=\frac{dV}{dx}\frac{1}{\sqrt{2}\,\sigma^{(1/4)}}\label{9}
\end{equation}

or,
\begin{equation}
\frac{dV}{dx}=\sqrt{2}\,\sigma^{(1/4)}\frac{dw}{dz}.\label{10}
\end{equation}

Therefore,
\begin{equation}
\frac{d^2w}{dz^2}=\frac{d}{dz}\left(\frac{dw}{dz}\right)=\frac{d}{dx}\left(\frac{dV}{dx}\frac{1}{\sqrt{2}\,\sigma^{(1/4)}}\right)\frac{dx}{dz}.\label{11}
\end{equation}

Simplifying Eq.\eqref{11} gives,
\begin{equation}
\frac{d^2 V}{dx^2}=2\sqrt{\sigma}\,\frac{d^2 w}{dz^2}\label{12}.
\end{equation}

Combining the above results, we have that
\begin{equation}
\frac{d^2 V}{dx^2}+(\sigma x^2 - \lambda)V= 2\sqrt{\sigma}\frac{d^2 w}{dz^2}+\left(\sigma \frac{z^2}{2\sqrt{\sigma}}-\lambda\right)\!w(z)=0,\label{13}
\end{equation}

therefore
\begin{equation}
\frac{d^2 w}{dz^2}+\left(\frac{z^2}{4}-\frac{\lambda}{2\sqrt{\sigma}}\right)\!w(z)=0.\label{14}
\end{equation}

Thus if $V$ is a solution of Eq. \eqref{6}, then $w(z):= V\!\left(z/(\sqrt{2}\,\sigma^{(1/4)})\right)$ is a solution of Eq. \eqref{4} with $a=\lambda/(2\sqrt{\sigma})$.

In a similar way, one can go in the reverse direction.

Let $z$ be a function of $x$ given by
\begin{equation}
z(x)=\sqrt{2}\,\sigma^{(1/4)}x,\label{15}
\end{equation}

and define,
\begin{equation}
V(x)=w(z(x))=w(\sqrt{2}\,\sigma^{(1/4)}x).\label{16}
\end{equation}

Then by the chain rule,
\begin{equation}
\frac{dV}{dx}=\frac{dw}{dz}\frac{dz}{dx}=\sqrt{2}\,\sigma^{(1/4)}\frac{dw}{dz},\label{17}
\end{equation}

or,
\begin{equation}
\frac{dw}{dz}=\frac{1}{\sqrt{2}\,\sigma^{(1/4)}}\frac{dV}{dx}.\label{18}
\end{equation}

Similarly,
\begin{equation}
\frac{d^2V}{dx^2}=\sqrt{2}\,\sigma^{(1/4)}\frac{d}{dx}\left(\frac{dw}{dz}\right)=\sqrt{2}\,\sigma^{(1/4)}\frac{d}{dz}\left(\frac{dw}{dz}\right)\frac{dz}{dx}= 2\sqrt{\sigma}\,\frac{d^2w}{dz^2}.\label{19}
\end{equation}

Thus,
\begin{equation}
\frac{d^2w}{dz^2}=\frac{1}{2\sqrt{\sigma}}\frac{d^2V}{dx^2}.\label{20}
\end{equation}

Therefore,
\begin{equation}
\frac{d^2 w}{dz^2}+\left(\frac{z^2}{4}-a\right)w(z)=\frac{1}{2\sqrt{\sigma}}\frac{d^2V}{dx^2}+\left(\frac{2\sqrt{\sigma}}{4}x^2 -a\right)\!V=0.\label{21}
\end{equation}

Simplifying gives,
\begin{equation}
\frac{d^2 V}{dx^2}+ (\sigma x^2 - 2a\sqrt{\sigma})V=0.\label{22}
\end{equation}

Thus if $w$ is a solution of Eq. \eqref{3}, then $V(x)= w(z(x))=w(\sqrt{2}\sigma^{(1/4)}x)$ is a solution of Eq. \eqref{6} with $\lambda=2a\sqrt{\sigma}$.
\end{proof}

\end{document}
